% Part: propositional-logic
% Chapter: syntax-and-semantics
% Section: formulas

\documentclass[../../../include/open-logic-section]{subfiles}

\begin{document}

\olfileid[tr]{pl}{syn}{fml}

\olsection{Önerme Formülleri}

Önerme mantığında !!{formula}s, başlangıç ifadeleri olarak
\emph{!!{propositional variable}s}\iftag{prvFalse}{\iftag{prvTrue}{,}{ ve}
  önerme sabiti~$\lfalse$}{}\iftag{prvTrue}{\iftag{prvFalse}{ ve}{ ve}
  önerme sabiti~$\ltrue$}{} seçilip bunlara \emph{mantıksal bağlaçlar}
uygulanarak kurulur.

\begin{enumerate}
\item !!^a{denumerable} ve !!{propositional variable}s{}nden oluşan bir
  küme~$\PVar$:
  $\Obj p_0$, $\Obj p_1$, \dots
\tagitem{prvFalse}{!!{falsity} için önerme sabiti~$\lfalse$.}{}
\tagitem{prvTrue}{!!{truth} için önerme sabiti~$\ltrue$.}{}
\item Mantıksal bağlaçlar:
  \startycommalist
  \iftag{prvNot}{\ycomma $\lnot$ (değilleme)}{}%
  \iftag{prvAnd}{\ycomma $\land$ (ve bağlacı)}{}%
  \iftag{prvOr}{\ycomma $\lor$ (veya bağlacı)}{}%
  \iftag{prvIf}{\ycomma $\lif$ (!!{conditional})}{}%
  \iftag{prvIff}{\ycomma $\liff$ (!!{biconditional})}{}%
\item Noktalama işaretleri: (, ) ve virgül.
\end{enumerate}

Önerme mantığının bu dilini $\Lang L_0$ ile gösteririz.

\iftag{defNot,defOr,defAnd,defIf,defIff,defTrue,defFalse,defEx,defAll}{%
Yukarıda sunulan ilkel bağlaçlara ek olarak aşağıdaki 
\emph{tanımlı} simgeleri de kullanırız:
\startycommalist
  \iftag{defNot}{\ycomma $\lnot$ (değilleme)}{}%
  \iftag{defAnd}{\ycomma $\land$ (ve bağlacı)}{}%
  \iftag{defOr}{\ycomma $\lor$ (veya bağlacı)}{}%
  \iftag{defIf}{\ycomma $\lif$ (!!{conditional})}{}%
  \iftag{defIff}{\ycomma $\liff$ (!!{biconditional})}{}%
  \iftag{defFalse}{\ycomma $\lfalse$ (!!{falsity})}{}%
  \iftag{defTrue}{\ycomma $\ltrue$ (!!{truth})}}{}.

\begin{tagblock}{defNot,defOr,defAnd,defIf,defIff,defTrue,defFalse,defEx,defAll}
\begin{explain}
Tanımlı bir simge resmî olarak dilin parçası değildir; biçimsel olmayan bir
kısaltma olarak sunulur. Böylece yalnızca ilkel simgeler kullanıldığında epey
uzayacak formülleri kısaltabiliriz. Bu açıkça bir üstünlüktür. Ancak daha
büyük üstünlük, kanıtların kısalmasıdır. Bir simge ilkelse kanıtlarda ayrı
olarak ele alınmalıdır. Dolayısıyla ilkel simgelerin sayısı arttıkça
kanıtlarımız da uzar.
\end{explain}
\end{tagblock}

% Alternate symbols

\begin{intro}
Yukarıda kullandıklarımızdan farklı terim ve simgelere aşina olabilirsiniz.
Mantık metinleri (ve öğretmenleri) ``değilleme'' için genellikle $\sim$,
  $\neg$ ve~!; ``ve bağlacı'' içinse $\wedge$, $\cdot$ ve $\&$ simgelerinden
birini kullanır. ``Koşul'' ya da ``içerme'' için yaygın olarak kullanılan
simgeler $\rightarrow$, $\Rightarrow$ ve $\supset$'dır.
\iftag{prvIff,defIff}{``İki yönlü koşul'', ``iki yönlü içerme'' veya
  ``(maddi) eşdeğerlik'' için kullanılan simgeler $\leftrightarrow$,
  $\Leftrightarrow$ ve $\equiv$'dir.}{}
\iftag{prvFalse,defFalse}{$\lfalse$ simgesi ``yanlışlık'', ``falsum'',
  ``saçmalık'' ya da ``alt'' diye de adlandırılır.}{}
\iftag{prvTrue,defTrue}{$\ltrue$ simgesi ``doğruluk'', ``verum'' ya da
  ``üst'' diye de adlandırılır.}{}
\end{intro}

\begin{defn}[Formül]
\ollabel{defn:formulas}
Önerme mantığının \emph{!!{formula}s} kümesi~$\Frm[L_0]$ tümevarımlı
olarak şöyle tanımlanır:
\begin{enumerate}
\tagitem{prvFalse}{$\lfalse$, atomik bir !!{formula} olur.}{}

\tagitem{prvTrue}{$\ltrue$, atomik bir !!{formula} olur.}{}

\item Her !!{propositional variable}~$\Obj p_i$ atomik bir
  !!{formula} sayılır.

\tagitem{prvNot}{$!A$ !!a{formula} ise $\lnot !A$ da
  !!a{formula} olur.}{}

\tagitem{prvAnd}{$!A$ ve $!B$ !!{formula}s ise $(!A \land
  !B)$ de !!a{formula} olur.}{}

\tagitem{prvOr}{$!A$ ve $!B$ !!{formula}s ise $(!A \lor !B)$
  de !!a{formula} olur.}{}

\tagitem{prvIf}{$!A$ ve $!B$ !!{formula}s ise $(!A \lif !B)$
  de !!a{formula} olur.}{}

\tagitem{prvIff}{$!A$ ve $!B$ !!{formula}s ise $(!A \liff !B)$
  de !!a{formula} olur.}{}

\tagitem{limitClause}{Başka hiçbir şey !!a{formula} değildir.}{}
\end{enumerate}
\end{defn}

\begin{explain}
!!{formula}s için verdiğimiz tanım bir
\emph{tümevarımlı tanımdır}. Özünde !!{formula}s kümesini sonsuz sayıda
aşamada kurarız. Başlangıç aşamasında bütün atomik formüllerin formül
olduğunu bildiririz; bu, tanımın ilk birkaç durumuna, yani
\iftag{prvTrue}{$\ltrue$, }{}%
\iftag{prvFalse}{$\lfalse$, }{}%
$\Obj p_i$ durumlarına karşılık gelir. Dolayısıyla ``atomik !!{formula}'',
bu biçimlerden herhangi birindeki !!{formula} demektir.

Tanımın öteki durumları, daha önce kurulmuş !!{formula}s kullanarak yeni
!!{formula}s kurma kurallarını verir. İkinci aşamada bu kuralları atomik
!!{formula}s kullanarak !!{formula}s kurmak için kullanabiliriz. Üçüncü aşamada
atomik formüllerden ve ikinci aşamada elde edilenlerden yeni formüller
kurarız; bu böyle sürer. Bir aşamada sonunda kurulmuş olan her şey ve yalnızca
bu şeyler !!{formula} kapsamındadır.
\end{explain}

İki yerli bir bağlaç~$\ast$ kullanılarak $!B$ ve $!C$'den kurulan
$(!B \ast !C)$ formülünü yazarken en dıştaki parantez çiftini çoğu zaman
atar ve yalnızca~$!B \ast !C$ yazarız.

\begin{tagblock}{defNot,defOr,defAnd,defIf,defIff,defTrue,defFalse,defEx,defAll}
\begin{defn}
Tanımlı işleçler kullanılarak kurulan formüller şöyle anlaşılmalıdır:

\begin{tagenumerate}{defTrue,defFalse,defNot,defOr,defAnd,defIf,defIff,defEx,defAll}

\tagitem{defTrue}{$\ltrue$,
  \iftag{prvFalse}{$\lnot\lfalse$'ın kısaltmasıdır.}{burada $!A$
    sabitlenmiş atomik !!{formula} olmak üzere, $(!A \lor \lnot !A)$'nın
    kısaltmasıdır.}}{}

\tagitem{defFalse}{$\lfalse$,
  \iftag{prvTrue}{$\lnot\ltrue$'ın kısaltmasıdır.}{burada $!A$
    sabitlenmiş atomik !!{formula} olmak üzere, $(!A \land \lnot !A)$'nın
    kısaltmasıdır.}}{}

\tagitem{defNot}{$\lnot !A$, $!A \lif \lfalse$'ın kısaltmasıdır.}{}

\tagitem{defOr}{$!A \lor !B$,
  \iftag{prvAnd}{$\lnot(\lnot !A \land \lnot !B)$'nın}{$\lnot !A \lif
    !B$'nın} kısaltmasıdır.}{}

\tagitem{defAnd}{$!A \land !B$,
  \iftag{prvOr}{$\lnot(\lnot !A \lor \lnot !B)$'nın}{$\lnot (!A \lif
    \lnot !B)$'nın} kısaltmasıdır.}{}

\tagitem{defIf}{$!A \lif !B$,
  \iftag{prvOr}{$\lnot !A \lor !B$'nın}{$\lnot (!A \land \lnot !B)$'nın}
  kısaltmasıdır.}{}

\tagitem{defIff}{$!A \liff !B$, $(!A \lif !B) \land (!B
  \lif !A)$'nın kısaltmasıdır.}{}
\end{tagenumerate}
\end{defn}
\end{tagblock}

\begin{defn}[Sözdizimsel özdeşlik]
$\ident$ simgesi, simge dizileri arasındaki sözdizimsel özdeşliği ifade
eder; yani $!A \ident !B$ ancak ve ancak $!A$ ile $!B$ aynı uzunlukta ve her
konumda aynı simgeyi içeren simge dizileriyse geçerlidir.
\end{defn}

$\ident$ simgesinin iki yanında ardışık birleştirmeyle elde edilen diziler
bulunabilir. Örneğin $!A \ident (!B \lor !C)$ şu demektir: $!A$ simge dizisi,
bu sırayla bir açma parantezi, $!B$ dizisi, $\lor$ simgesi, $!C$ dizisi ve
bir kapama parantezinin ardışık birleştirilmesiyle elde edilen diziyle
aynıdır. Durum buysa $!A$'nın ilk simgesinin bir açma parantezi olduğunu,
$!A$'nın $!B$'yi bir alt dizi olarak (ikinci simgeden başlayarak) içerdiğini,
bu alt diziyi $\lor$'un izlediğini vb. biliriz.

\end{document}
